% 本书图片的版权说明（仅第 1 册）。
% 机器可读的授权记录位于 images/book1/CREDITS.md；增加或删除照片时，请保持两个文件同步。
\chapter*{图片来源}

本书中的插图均为原创。照片依其各自的自由许可使用，谨向作者致谢：

\begin{itemize}
  \item \emph{吸住回形针的磁铁}（第 6 章）：Donald
        Olszewski，Wikimedia Commons，CC~BY~4.0。
  \item \emph{沙漏}（第 8 章）：John Morgan，Wikimedia
        Commons，CC~BY~2.0。
  \item \emph{月球正面}（第 11 章）：NASA/GSFC/Arizona
        State University（Lunar Reconnaissance Orbiter），公有领域。
  \item \emph{带地照的蛾眉月}（第 11 章）：
        G.~Donatiello，Wikimedia Commons，CC0。
  \item \emph{白昼中的凸月}（第 11 章）：Cbaile19，
        Wikimedia Commons，CC0。
  \item \emph{溴：一种有颜色的气体}（第 14 章）：James
        St.~John，Wikimedia Commons，CC~BY~2.0。
  \item \emph{鼓和鼓槌}（第 19 章）：Vladimir Morozov，
        Wikimedia Commons，CC~BY~2.0。
  \item \emph{指南针}（第 25 章）：Evan-Amos，Wikimedia Commons，
        公有领域。
  \item \emph{太阳}（第 26 章）：NASA/SDO (AIA)，公有领域。
  \item \emph{八大行星}（第 26 章）：NASA——
        水星：NASA/JHU-APL/Carnegie（MESSENGER）；金星：
        NASA/JPL（Magellan 雷达拼接图）；地球：NASA，Apollo~17；
        火星：NASA/USGS（Viking）；木星：NASA/JPL/Space Science
        Institute（Cassini）；土星：NASA/JPL/Space Science
        Institute（Cassini）；天王星和海王星：NASA/JPL
        （Voyager~2）。均为公有领域。
\end{itemize}

每张照片的完整来源链接和许可证 URL 均列于本书代码库中的
\texttt{images/book1/CREDITS.md}。

\bigskip
除照片外，第 1--5、7、9--10、12--13、15--18、20--22、24、27--29、31--34、36--40、42、44--45、48、50、53--55、59--61、63、66 以及 69--70 章中的四十五幅图为 AI 生成的插图，由 OpenAI 图像生成功能根据本书编辑撰写的提示词生成，并在收录前已审核其物理准确性。每幅生成图像的完整提示词均记录于代码库中的
\texttt{images/book1/ai/PROMPTS.md}。
\clearpage
